% Generated draft: requires mathematical and source review.
\documentclass[11pt]{article}
\usepackage[margin=25mm]{geometry}
\usepackage{fontspec}
\setmainfont{Latin Modern Roman}
\usepackage{amsmath,amssymb,mathtools,unicode-math}
\setmathfont{Latin Modern Math}
\usepackage{enumitem,needspace,xurl,hyperref}
\setlength{\parindent}{0pt}
\setlength{\parskip}{5pt}
\setlength{\emergencystretch}{3em}
\newcommand{\N}{\mathbb{N}}
\newcommand{\Z}{\mathbb{Z}}
\newcommand{\Q}{\mathbb{Q}}
\newcommand{\R}{\mathbb{R}}
\newcommand{\C}{\mathbb{C}}
\newcommand{\F}{\mathbb{F}}
\newcommand{\A}{\mathbb{A}}
\newcommand{\PP}{\mathbb{P}}
\newcommand{\E}{\mathbb{E}}
\newcommand{\eps}{\varepsilon}
\newcommand{\dd}{\,\mathrm{d}}
\begin{document}
\section*{UnsolvedMath ID: 5300025}
\textbf{Source problem code:} catalogue entry
\subsection{Definitions and mathematical statement}
A domain \(U\) is John, in a specified metric \(d\), if there are \(z_0\in U\) and \(C\ge1\) such that each \(z\in U\) can be joined to \(z_0\) by a rectifiable curve \(\gamma:[0,L]\to U\), parametrized by arclength from \(z\), satisfying
\[
\operatorname{dist}(\gamma(s),\partial U)\ge s/C\qquad(0\le s\le L).
\]
The metric and the convention for an unbounded domain or a base point at infinity are part of this definition. For a rational map \(f\), its Fatou set \(F(f)\) is the open normality locus of its iterates and \(J(f)\) its Julia set, the complement. Individual Fatou components are domains, but a disconnected union is not a John domain in this sense; the Julia set itself is not the open domain in question.

For a specified critical point \(c_*\), the literal uniform nonrecurrence condition can be written
\[
\exists\delta>0\ \forall n\ge1,\qquad d(f^n(c_*),c_*)\ge\delta.
\]
In the usual quadratic parameter family \(f_c(z)=z^2+c\), this concerns \(c_*=0\). A strict postcritically finite Misiurewicz convention means that this point is strictly preperiodic to a repelling periodic orbit; a broader nonrecurrent-critical convention is different. The extraction does not say which convention or map family is intended.

Carleson's question asks how far a real-quadratic equivalence between critical nonrecurrence and a John property extends to general Misiurewicz parameters. To make an equivalence precise one must specify the map class and the domain \(U(f)\), and then ask whether
\[
c_*\text{ satisfies the displayed nonrecurrence condition}
\quad\Longleftrightarrow\quad U(f)\text{ is John}.
\]
The source does not identify whether \(U(f)\) is the basin of infinity, each Fatou component, or another chosen domain, nor supply exclusions such as parabolic behavior or rules for multiple critical points. These are not silently filled in.

The supplied research summary reports a solved real-quadratic case and partial expanding/nonrecurrent results, while calling the general equivalence unresolved. Its phrases Julia set is a John domain and Julia set is a quasicircle/John domain conflate different geometric objects; they are not adopted as precise theorems. Its reported partial results require their actual hypotheses and domain conventions.
\subsection{Short English statement}
Determine when nonrecurrence of critical points characterizes a John property for the intended Fatou domain beyond the real-quadratic setting. The parameter class, Misiurewicz convention, and chosen domain need clarification.
\subsection{Sources}
\begin{itemize}
\item Local UnsolvedMath registry, permanent ID \(5300025\), exact source code \texttt{catalogue entry}, ``John domains at general Misiurewicz points.''
\item Stony Brook, \emph{Open Problems in Dynamical Systems}, 1992, Carleson 4; proposer Lennart Carleson. \url{https://www.math.stonybrook.edu/open-problems-dynamical-systems}. Recorded access 2026-07-29; source-TeX-and-pdf-text extraction.
\item Supplied research summary mentions Zhang; Carleson--Jones--Yoccoz, ``Julia and John''; and Przytycki. No complete references or exact theorem hypotheses are supplied.
\end{itemize}
\end{document}
